%%%%%%%%%%%%%%%%%%%%%%%%%%%%%%%%%%%%%%%%%
% University/School Laboratory Report
% LaTeX Template
% Version 3.1 (25/3/14)
%
% This template has been downloaded from:
% http://www.LaTeXTemplates.com
%
% Original author:
% Linux and Unix Users Group at Virginia Tech Wiki 
% (https://vtluug.org/wiki/Example_LaTeX_chem_lab_report)
%
% License:
% CC BY-NC-SA 3.0 (http://creativecommons.org/licenses/by-nc-sa/3.0/)
%
%%%%%%%%%%%%%%%%%%%%%%%%%%%%%%%%%%%%%%%%%

%----------------------------------------------------------------------------------------
%	PACKAGES AND DOCUMENT CONFIGURATIONS
%----------------------------------------------------------------------------------------

\documentclass{article}

\usepackage[version=3]{mhchem} % Package for chemical equation typesetting
\usepackage{siunitx} % Provides the \SI{}{} and \si{} command for typesetting SI units
\usepackage{graphicx} % Required for the inclusion of images
\usepackage{amsmath} % Required for some math elements 
\usepackage{verbatim}
\usepackage{enumitem}

\setlength\parindent{0pt} % Removes all indentation from paragraphs

\renewcommand{\labelenumi}{\alph{enumi}.} % Make numbering in the enumerate environment by letter rather than number (e.g. section 6)

%\usepackage{times} % Uncomment to use the Times New Roman font

%----------------------------------------------------------------------------------------
%	DOCUMENT INFORMATION
%----------------------------------------------------------------------------------------

\title{Optional Lab Template} % Title

\author{Ryan \textsc{Rodriguez}} % Author name

\date{\today} % Date for the report

\begin{document}

\maketitle % Insert the title, author and date

\begin{center}
\begin{tabular}{l r}
Date Performed: & January 1, 2012 \\ % Date the experiment was performed
TA: & Ryan Rodriguez % Instructor/supervisor
\end{tabular}
\end{center}

% If you wish to include an abstract, uncomment the lines below
% \begin{abstract}
% Abstract text
% \end{abstract}

%----------------------------------------------------------------------------------------
%	Purpose: Briefly describe what the desired learning outcomes for the lab assigment are.
%	Learning outcomes will often becomes clear by simply reading the lab assignment first.
%----------------------------------------------------------------------------------------

\section{Purpose}

The purpose of this lab is to \underline{\hspace{3cm}}.  We will learn skills $\left \{skill_1, \ldots, skill_n \right \}$ by doing the following steps \underline{\hspace{3cm}} described in the lab manual.

 
%----------------------------------------------------------------------------------------
%	SECTION 2
%----------------------------------------------------------------------------------------

\section{Experiments}
Here, you should do the following:
\begin{itemize}
     \item Create a section for each major section in the lab - you can copy and paste the section you are reading now.
     \item In each of these new appropriately titled sections:
     \begin{itemize}
     	\item fill in some of the major steps you took to complete the exercise 
     	\item Feel free to discuss any challenges you faced and overcame in completing the lab.
        \item Include images when they are helpful to illustrate your point. Schematics from MML might be something you'd like to include. 
     \end{itemize}
    
    \item Finally, create a subsection titled 'results' for each section you've created.
	\begin{itemize}
    	\item In the results section, describe the implications of the work you've done.
        \item \emph{Note:} You can create subsection by doing \verb|\subsection{TITLE}| instead of \verb|\section{TITLE}|.
	\end{itemize}
\end{itemize}

\subsection{Results}
For each experiment, add a results section.


%----------------------------------------------------------------------------------------
%	SECTION 3
%----------------------------------------------------------------------------------------

\section{Conclusion}

Sum up what you did, what you learned.

%----------------------------------------------------------------------------------------
%	Answers to Questions
%----------------------------------------------------------------------------------------

\section{Answers to Questions}

Answer any questions, if applicable. You can do so, for example, like this:

\begin{description}[style=nextline]
    \item[Question 1: What is Euler's formula?] Answer: $e^{ \pm i\theta } = \cos \theta \pm i\sin \theta$
    \item[Question 2: How many bits in a nibble?] Answer: 4
    \item[Question 3: Who is your favorite TA?] Answer:  Ryan Rodriguez
    
    \end{description}

%----------------------------------------------------------------------------------------
%	BIBLIOGRAPHY
%----------------------------------------------------------------------------------------

 \begin{thebibliography}{1}
  
  \bibitem{notes} Albert Camus {\em The Rebel: An Essay on Man in Revolt.}  New York: Vintage, 1956. Print.

  \end{thebibliography}
  

%----------------------------------------------------------------------------------------


\end{document}